\pdfvariable trailerid {[<C2532026083000000000000000000000><C2532026083000000000000000000000>]}
\documentclass[10pt]{article}
\usepackage[margin=0.72in]{geometry}
\usepackage{amsmath,amssymb,amsthm,booktabs,hyperref}
\hypersetup{colorlinks=true,linkcolor=blue,urlcolor=blue}
\newtheorem{theorem}{Theorem}
\title{Exact Fixation and Green Kernels for a Finite Moran Process}
\author{Route-A source-local certificate HCS-C253}
\date{30 August 2026}
\begin{document}
\maketitle
\begin{abstract}
We close an exact finite-state atlas for a continuous-time Moran birth--death
process with selection ratio $\rho$ and event scale $\beta$. Fixation,
occupation, absorption time, and reversible weights are rational for rational
inputs. This is source-local probability; no arithmetic origin or spectral
operator is claimed.
\end{abstract}
\section{Model}
Let $X_t\in\{0,\ldots,N\}$ count one type in a fixed population. For
$1\le i\le N-1$, set
\[
\lambda_i=\beta\rho\,\frac{i(N-i)}N,\qquad
\mu_i=\beta\,\frac{i(N-i)}N. \tag{1}
\]
The endpoints $0$ and $N$ absorb. The transient generator $Q$ has
$Q_{i,i+1}=\lambda_i$, $Q_{i,i-1}=\mu_i$, and
$Q_{ii}=-(\lambda_i+\mu_i)$.
\section{Fixation theorem}
Let $u_i$ be the probability of absorption at $N$. The backward equation is
\[
\lambda_i(u_{i+1}-u_i)+\mu_i(u_{i-1}-u_i)=0,\quad u_0=0,\ u_N=1.
\]
\begin{theorem}[selection-ratio fixation]
For $\rho\ne1$, $u_i=(1-\rho^{-i})/(1-\rho^{-N})$; for $\rho=1$,
$u_i=i/N$.
\end{theorem}
\begin{proof}
Successive differences have ratio $\rho^{-1}$; summing the geometric
progression and taking the neutral limit gives the formula. Positivity of
all transient rates and finiteness of the state space imply absorption almost
surely.
\end{proof}
\section{Green kernel and time}
Define $G=(-Q)^{-1}$ and $t_i=\sum_{j=1}^{N-1}G_{ij}$. Then $G_{ij}$ is
expected occupation time and
\[
\lambda_i(t_{i+1}-t_i)+\mu_i(t_{i-1}-t_i)=-1,\quad t_0=t_N=0. \tag{2}
\]
All entries are rational for rational $(N,\rho,\beta)$.
The factorization separates population selection from the arbitrary event-time
scale, so the Green kernel is not a fitted schedule.
\section{Reversibility and boundaries}
With $w_1=1$ and
\[
\frac{w_{i+1}}{w_i}=\rho\,\frac{i(N-i)}{(i+1)(N-i-1)}, \tag{3}
\]
one has $w_i\lambda_i=w_{i+1}\mu_{i+1}$. The $\rho=1$ face is neutral;
$\beta=0$ freezes the chain, and $N=1$ has no transient block.
\begin{table}[htbp]
\centering\small
\caption{Receipt coverage.}
\begin{tabular}{llll}
\toprule
rows & sizes & selection & recorded object\\
\midrule
2 & 3,8 & neutral & fixation, Green, time\\
4 & 4,6,7,10 & mixed & fixation, full Green\\
2 & 5,9 & weak/strong & fixation, full Green\\
\bottomrule
\end{tabular}
\end{table}
\section{Executable certificate}
The producer emits eight exact rows, including every Green matrix through
$N=10$. The independent checker closes 220 assertions; SymPy closes ten
identities; replay is byte-identical and hostile mutation rejects 23/23.
Three fixed-epoch manuscript rounds use embedded fonts and a content-addressed
manifest. Decimal values are presentation checks only. Every rational
recurrence is checked before decimal conversion, including the neutral limit
and the no-inversion boundary faces.
\section{Route-A boundary}
The strict evaluator tuple is
\texttt{(A0\_FAIL,A1\_PASS\_ANALYTIC,A2\_FAIL,A3\_FAIL,A4\_FORMAL\_HINT)}.
The arithmetic origin is \texttt{none}; no primitive target orbit or
determinant convention is defined. The verdict is
\texttt{ROUTE\_A\_REJECTED}, with \texttt{route\_b\_invocation\_allowed: false},
under \texttt{NO\_BAD\_EULER\_OR\_ROOT\_NUMBER}. Fixation is a population
probability, not a target arithmetic matching law.
\section{Conclusion}
The finite Moran generator supplies both a closed selection formula and an
exact occupation kernel. Explicit frozen and singleton faces preserve the A0
stopping boundary.
\begin{thebibliography}{9}
\bibitem{Moran} P. A. P. Moran, ``Random processes in genetics,'' Math.
Proc. Cambridge Philos. Soc. 54 (1958).
\bibitem{Ewens} W. J. Ewens, \emph{Mathematical Population Genetics},
Springer, 2004.
\bibitem{Norris} J. R. Norris, \emph{Markov Chains}, Cambridge Univ. Press,
1997.
\end{thebibliography}
\end{document}
